\begin{exercice*}
    Montre que les deux solides ci-dessous ont le même volume.
    \begin{center}
        \Solide[%
            Nom=pave,
            Phi=110,
            Largeur=1.3,
            Hauteur=0.975,
            Profondeur=1.95,
            Sommets=false,
            Traces={
                    trace appelation(Projette(C),Projette(D),-3mm,TEX("$2x$"));
                    trace appelation(Projette(B),Projette(C),-3mm,TEX("$2x+4$"));
                    trace appelation(Projette(D),Projette(E),-3mm,TEX("\rotatebox{90}{$x$}"));
                    Label(TEX("$\mathbf{\mathcal{V}_1}$"),iso(E,F,G,H));
                }
        ]

        \begin{Geometrie}[Cadre="aucun",TypeTrace="Espace"]
            Initialisation(5,45,20,57);
            color A,B,C,D,E,F;
            B=(0,0,0);
            A-B=(0,0,1);
            C-B=(4/5,0,0);
            E-B=(0,9/5,0);
            D-E=A-B;
            F-E=C-B;
            NbS:=6;
            Sommet1:=A;
            Sommet2:=B;
            Sommet3:=C;
            Sommet4:=D;
            Sommet5:=E;
            Sommet6:=F;
            NF:=5;
            Fc[100]:=3;Fc[101]:=1;Fc[102]:=2;Fc[103]:=3;
            Fc[200]:=3;Fc[201]:=6;Fc[202]:=5;Fc[203]:=4;
            Fc[300]:=4;Fc[301]:=4;Fc[302]:=5;Fc[303]:=2;Fc[304]:=1;
            Fc[400]:=4;Fc[401]:=5;Fc[402]:=6;Fc[403]:=3;Fc[404]:=2;
            Fc[500]:=4;Fc[501]:=1;Fc[502]:=3;Fc[503]:=6;Fc[504]:=4;
            DessineObjet;
            trace codeperp(A,B,C,8);
            trace codeperp(A,B,E,8);
            trace codeperp(D,E,F,8);
            label.rt(TEX("$4x$"),Projette(iso(D,E)));
            trace appelation(A,D,2mm,TEX("$x+2$"));
            trace appelation(F,E,-3mm,TEX("$2x$"));
            label(TEX("$\mathbf{\mathcal{V}_2}$"),Projette(iso(A,B,E,D)));
            %   label.top(btex A etex,Projette(A));
            %   label.top(btex D etex,Projette(D));
            %   label.rt(btex E etex,Projette(E));
            %   label.lrt(btex F etex,Projette(F));
            %   label.llft(btex C etex,Projette(C));
            %   label.bot(btex B etex,Projette(B));
        \end{Geometrie}
    \end{center}

\end{exercice*}
\begin{corrige}
    %\setcounter{partie}{0} % Pour s'assurer que le compteur de \partie est à zéro dans les corrigés
    %\phantom{rrr}    
    \dots
\end{corrige}

